\documentclass[10pt, aspectratio=169]{beamer}
\usepackage{../beamer_theme_408}

\title{408考研零基础 --- 计算机组成原理}
\subtitle{2-7 主存储器}
\author{}
\date{}

\begin{document}


\begin{frame}{本节目标}
    \begin{enumerate}
        \item 掌握SRAM与DRAM的区别
        \item 理解存储芯片的位扩展与字扩展
        \item 掌握主存储器与CPU的连接方法
        \item 了解内存条（SIMM/DIMM）的基本概念
    \end{enumerate}
\end{frame}

\begin{frame}[t]{SRAM的存储原理}
    \begin{block}{基本存储元——6管触发器}
        \begin{itemize}
            \item 由6个MOS管构成RS触发器
            \item 两个交叉耦合的反相器保持0或1状态
            \item 信息稳定存储在正反馈环路中
        \end{itemize}
    \end{block}
    \vspace{0.3em}
    \begin{block}{特点}
        \begin{itemize}
            \item \textcolor{red}{不需要刷新}——静态存储
            \item 存取速度快（纳秒级）
            \item 集成度低（6管/位），功耗较大
            \item 价格昂贵——主要用于Cache
        \end{itemize}
    \end{block}
\end{frame}

\begin{frame}[t]{DRAM的存储原理}
    \begin{block}{基本存储元——1管1电容}
        \begin{itemize}
            \item 由一个MOS管和一个电容组成
            \item 电容充电代表1，放电代表0
            \item 电容会漏电——需要定期刷新
        \end{itemize}
    \end{block}
    \vspace{0.3em}
    \begin{block}{刷新方式}
        \begin{itemize}
            \item 集中刷新：连续刷完所有行，期间不能访问
            \item 分散刷新：每读一次就刷新一次
            \item 异步刷新：在固定间隔刷新一行，轮流进行
        \end{itemize}
    \end{block}
    \vspace{0.3em}
    \begin{itemize}
        \item 典型刷新周期：64ms内刷新所有行
    \end{itemize}
\end{frame}

\begin{frame}{SRAM vs DRAM 总结}
    \begin{center}
        \begin{tabular}{|c|c|c|}
            \hline
            \textbf{对比项} & \textbf{SRAM} & \textbf{DRAM} \\
            \hline
            存储元 & 6管触发器 & 1管+电容 \\
            \hline
            是否刷新 & 不需要 & 需要（64ms） \\
            \hline
            存取速度 & 快（1$\sim$10ns） & 较慢（50$\sim$100ns） \\
            \hline
            集成度 & 低 & 高 \\
            \hline
            功耗 & 大 & 小 \\
            \hline
            价格 & 高 & 低 \\
            \hline
            用途 & Cache & 主存 \\
            \hline
        \end{tabular}
    \end{center}
\end{frame}

\begin{frame}{存储芯片的基本结构}
    \begin{center}
        \begin{tabular}{|c|}
            \hline
            地址线 $A_0 \sim A_{k-1}$ \\
            数据线 $D_0 \sim D_{n-1}$ \\
            片选信号 $\overline{CS}$ \\
            读写控制 $\overline{WE}$ \\
            \hline
        \end{tabular}
    \end{center}
    \vspace{0.5em}
    \begin{itemize}
        \item 地址线 $k$ 条 $\Rightarrow$ $2^k$ 个存储单元
        \item 数据线 $n$ 条 $\Rightarrow$ 每单元 $n$ 位
        \item 芯片容量 = $2^k \times n$ 位
    \end{itemize}
    \vspace{0.3em}
    \begin{exampleblock}{例：$2K \times 8$位芯片}
        \begin{align*}
            2K &= 2^{11} \Rightarrow 11\text{条地址线} \\
            8 &\Rightarrow 8\text{条数据线}
        \end{align*}
    \end{exampleblock}
\end{frame}

\begin{frame}{位扩展}
    \begin{block}{为什么要位扩展？}
        \begin{itemize}
            \item 芯片数据线位数不够时，用多片拼出所需字长
        \end{itemize}
    \end{block}
    \vspace{0.3em}
    \begin{exampleblock}{例：用 $1K \times 4$ 位芯片构成 $1K \times 8$ 位存储器}
        \begin{itemize}
            \item 需要芯片数 = $\dfrac{1K \times 8}{1K \times 4} = 2$ 片
            \item 连接方式：
            \begin{itemize}
                \item 两片地址线\textbf{并联}（共用地址）
                \item 片选信号\textbf{并联}（同时选中）
                \item 数据线各接高4位和低4位
            \end{itemize}
        \end{itemize}
    \end{exampleblock}
    \vspace{0.3em}
    \begin{center}
        \textcolor{blue}{位扩展：地址线并联，数据线分开，同时工作}
    \end{center}
\end{frame}

\begin{frame}{字扩展}
    \begin{block}{为什么要字扩展？}
        \begin{itemize}
            \item 芯片容量（地址空间）不够时，用多片增加地址范围
        \end{itemize}
    \end{block}
    \vspace{0.3em}
    \begin{exampleblock}{例：用 $1K \times 8$ 位芯片构成 $4K \times 8$ 位存储器}
        \begin{itemize}
            \item 需要芯片数 = $\dfrac{4K \times 8}{1K \times 8} = 4$ 片
            \item 连接方式：
            \begin{itemize}
                \item 各片数据线\textbf{并联}
                \item 地址线通过\textbf{片选译码}区分
                \item 高位地址线接译码器，产生片选信号
            \end{itemize}
            \item 地址分配：
            \begin{itemize}
                \item 芯片0：$0 \sim 1023$
                \item 芯片1：$1024 \sim 2047$
                \item 芯片2：$2048 \sim 3071$
                \item 芯片3：$3072 \sim 4095$
            \end{itemize}
        \end{itemize}
    \end{exampleblock}
\end{frame}

\begin{frame}{同时进行位扩展与字扩展}
    \begin{exampleblock}{例：用 $1K \times 4$ 位芯片构成 $4K \times 8$ 位存储器}
        \begin{itemize}
            \item 总位数：$4K \times 8 = 32K$ 位
            \item 每片位数：$1K \times 4 = 4K$ 位
            \item 需要芯片数 = $32K / 4K = 8$ 片
        \end{itemize}
    \end{exampleblock}
    \vspace{0.3em}
    \begin{center}
        \begin{tabular}{ccl}
            分组： & & 每2片做\textbf{位扩展} $\to$ 1组（$1K \times 8$） \\
                  & & 共4组做\textbf{字扩展} $\to$ $4K \times 8$ \\
        \end{tabular}
    \end{center}
    \vspace{0.3em}
    \begin{itemize}
        \item 组内：地址线并联，片选并联，数据线分高低
        \item 组间：地址译码片选
    \end{itemize}
\end{frame}

\begin{frame}{主存储器与CPU的连接}
    \begin{block}{需要考虑的问题}
        \begin{enumerate}
            \item \textbf{地址线连接}：CPU地址线 $\to$ 芯片地址线
            \item \textbf{数据线连接}：CPU数据线 $\to$ 芯片数据线
            \item \textbf{片选译码}：高位地址 $\to$ 译码器 $\to$ 片选
            \item \textbf{读写控制}：CPU读/写信号 $\to$ 芯片$\overline{WE}$
        \end{enumerate}
    \end{block}
    \vspace{0.3em}
    \begin{block}{片选译码方式}
        \begin{itemize}
            \item \textbf{线选法}：用一根高位地址线直接作为片选
            \begin{itemize}
                \item 简单但地址空间不连续，有重叠
            \end{itemize}
            \item \textbf{译码法}：用译码器（如3-8译码器）
            \begin{itemize}
                \item 地址连续，无地址重叠
            \end{itemize}
        \end{itemize}
    \end{block}
\end{frame}

\begin{frame}{内存条（SIMM/DIMM）}
    \begin{block}{SIMM (Single Inline Memory Module)}
        \begin{itemize}
            \item 引脚分布在模块的单侧
            \item 数据宽度：32位（72线）
            \item 常用于Pentium之前的计算机
        \end{itemize}
    \end{block}
    \begin{block}{DIMM (Dual Inline Memory Module)}
        \begin{itemize}
            \item 引脚分布在模块的两侧（电气独立）
            \item 数据宽度：64位（168线以上）
            \item SDRAM时代开始使用，至今仍是主流
        \end{itemize}
    \end{block}
    \vspace{0.3em}
    \begin{center}
        \begin{tabular}{|c|c|c|}
            \hline
            \textbf{类型} & \textbf{引脚} & \textbf{数据宽度} \\
            \hline
            SIMM & 72 & 32位 \\
            \hline
            DIMM & 168+ & 64位 \\
            \hline
        \end{tabular}
    \end{center}
\end{frame}

\begin{frame}{本节总结}
    \begin{enumerate}
        \item SRAM（6管触发器）vs DRAM（1管+电容，需刷新）
        \item 位扩展（增位数）和字扩展（增容量）
        \item CPU连接要点：地址线、数据线、片选译码、读写控制
        \item SIMM（单侧引脚）和DIMM（双侧引脚）
    \end{enumerate}
\end{frame}

\end{document}
